%% =========================================================================
%% Standalone abstract for conference submission.
%% Compile with: pdflatex abstract_submission.tex
%% This is self-contained — no external stats files needed.
%% =========================================================================

\documentclass[11pt]{article}
\usepackage[margin=1in]{geometry}
\usepackage{parskip}
\usepackage[hidelinks]{hyperref}

\title{Cyclical Demand Shocks and Aggregate Task Composition:\\
       Evidence from a Bartik Instrument}

\author{Ethan Johnston\thanks{%
  Department of Economics, Illinois State University, Campus Box 4200,
  Normal, IL 61790, United States. Email: \href{mailto:egjohn3@ilstu.edu}{\texttt{egjohn3@ilstu.edu}}.
  ORCID: \href{https://orcid.org/0009-0000-7599-5346}{0009-0000-7599-5346}.}}
\date{}

\begin{document}
\maketitle
\thispagestyle{empty}

\begin{abstract}
\noindent I estimate the effect of cyclical labor demand shocks on the task composition of the employed workforce using local projections instrumented with a Bartik shift-share predictor. Exploiting variation in pre-period industry composition across 51 states from 1978 to 2025, I find that a one-percentage-point increase in state unemployment reduces the mean routine task content of employed workers by $-0.0091$ at impact (first-stage $F = 23.3$), with the effect persisting and modestly intensifying over a five-year horizon rather than reverting. The aggregate shift is concentrated in the routine-manual employment share, which falls sharply during downturns and remains depressed; the decline is offset by rising non-routine employment, with the non-routine-manual share the most persistent absorbing margin. Controlling for cumulative changes in industry employment composition attenuates the impact effect by only 12.9\%, indicating that the response operates primarily within rather than between industries. Placebo tests reveal no significant differential pre-trend in routine task content, the labor-force-participation response is too small at all horizons to drive the headline pattern, and the result survives flexible technology controls, state-specific trends, region-by-year trends, leave-one-industry-out variants of the Bartik instrument, a drop-COVID subsample, and weak-instrument-robust inference. The findings are consistent with cyclical demand shocks contributing to long-run task reallocation rather than generating transitory deviations around a structural trend.
\end{abstract}

\medskip
\noindent\textbf{Keywords:} Task composition; Routine employment; Business cycles; Local projections; Bartik instrument; Labor reallocation.

\medskip
\noindent\textbf{JEL Codes:} E24, E32, J23, J24, J62.

\end{document}
